%% ============================================================
%% preamble/theorems.tex — theorem environments (loaded LAST, after cleveref)
%%
%% Purpose      : the amsthm environments this paper uses, numbered in one
%%                shared AMS series, and named correctly by \cref/\Cref.
%% Dependencies : amsthm (packages.tex), cleveref (hyperref-last.tex),
%%                aliascnt (loaded here).
%%
%% TWO THINGS ARE LOAD-BEARING HERE, AND BOTH WERE MEASURED, NOT ASSUMED.
%%
%% 1. This file is input AFTER cleveref (see preamble/main.tex). cleveref
%%    patches \newtheorem when it loads; an environment declared earlier is
%%    invisible to that patch.
%%
%% 2. That alone is NOT enough, which is the part that is easy to get
%%    wrong. cleveref names a reference after the COUNTER the environment
%%    uses, and the plain AMS idiom
%%
%%        \newtheorem{lemma}[theorem]{Lemma}
%%
%%    makes every environment share the `theorem' counter — which is
%%    exactly what we want for NUMBERING (so that Lemma 2.1 and Theorem 2.2
%%    sit in one readable series) and exactly wrong for NAMING. The first
%%    build of this paper printed "Theorem 2.1" for \Cref{lem:support} and
%%    "Theorem 3.4" for a remark, in the middle of a proof, and moving this
%%    file after cleveref did not change it.
%%
%%    \newaliascnt is the fix: it gives each environment its OWN counter
%%    that is an alias of `theorem', so the numbering is still shared and
%%    cleveref sees a distinct type per environment. \aliascntresetthe is
%%    required after each alias (aliascnt manual) or \thelemma expands
%%    wrongly. The \crefname lines below then bind the printed names.
%%
%% Check this file by reading a \Cref to a lemma in the built PDF, never by
%% reading the source: the failure mode is a correct-looking build.
%% ============================================================

\usepackage{aliascnt}

% ---------------------------------------------------------------------
% The master counter, numbered within the section.
% ---------------------------------------------------------------------
\theoremstyle{plain}
\newtheorem{theorem}{Theorem}[section]

% Aliased counters: shared numbering, distinct cleveref types.
\newaliascnt{lemma}{theorem}
\newtheorem{lemma}[lemma]{Lemma}
\aliascntresetthe{lemma}

\newaliascnt{proposition}{theorem}
\newtheorem{proposition}[proposition]{Proposition}
\aliascntresetthe{proposition}

\newaliascnt{corollary}{theorem}
\newtheorem{corollary}[corollary]{Corollary}
\aliascntresetthe{corollary}

\theoremstyle{definition}
\newaliascnt{definition}{theorem}
\newtheorem{definition}[definition]{Definition}
\aliascntresetthe{definition}

\newtheorem{assumption}{Assumption}

\theoremstyle{remark}
\newaliascnt{remark}{theorem}
\newtheorem{remark}[remark]{Remark}
\aliascntresetthe{remark}

% ---------------------------------------------------------------------
% Printed names. `capitalize' (hyperref-last.tex) makes \Cref capitalise at
% a sentence start; the lower-case forms below are what \cref uses mid
% sentence.
% ---------------------------------------------------------------------
\crefname{theorem}{Theorem}{Theorems}
\Crefname{theorem}{Theorem}{Theorems}
\crefname{lemma}{Lemma}{Lemmas}
\Crefname{lemma}{Lemma}{Lemmas}
\crefname{proposition}{Proposition}{Propositions}
\Crefname{proposition}{Proposition}{Propositions}
\crefname{corollary}{Corollary}{Corollaries}
\Crefname{corollary}{Corollary}{Corollaries}
\crefname{definition}{Definition}{Definitions}
\Crefname{definition}{Definition}{Definitions}
\crefname{remark}{Remark}{Remarks}
\Crefname{remark}{Remark}{Remarks}
\crefname{assumption}{Assumption}{Assumptions}
\Crefname{assumption}{Assumption}{Assumptions}
